% !TeX root = main.tex
\changecolours{ScoutGreen}
\section[Visão Geral da Revisão N1]{Visão Geral da Revisão Integrada N1}

% ------------------------------------------------------------------------------
\begin{frame}{Objetivos do Sábado Letivo: Revisão e Simulado}
  \vspace{0.1cm}
  \begin{columns}[t,onlytextwidth]
    \begin{column}{0.48\textwidth}
      \alert{Propósito do Encontro}\par
      \footnotesize
      \begin{itemize}
        \item Consolidar o conteúdo integral do 1º Período Letivo (Aulas 01 a 08).
        \item \textbf{Sem conteúdo novo:} foco total em resolução prática e fixação de conceitos.
        \item Preparação imediata para os exames:
          \begin{itemize}
            \item \textbf{AVP1 (Aula 10):} Avaliação prática em laboratório com Redis e MongoDB.
            \item \textbf{AVT1 (Aula 11):} Avaliação teórica individual e entrega da Lista 1.
          \end{itemize}
      \end{itemize}
    \end{column}
    \begin{column}{0.48\textwidth}
      \alert{Matriz de Competências de N1}\par
      \footnotesize
      \begin{itemize}
        \item \textbf{Teoria:} Teorema CAP, Teorema PACELC e Propriedades ACID vs BASE.
        \item \textbf{Redis:} Chave-Valor, TTL, Hashes, Filas e Padrão Cache-Aside.
        \item \textbf{MongoDB:} BSON, Atualizações Atômicas, Regra ESR e Covered Queries.
        \item \textbf{Analytics:} Aggregation Framework e Pipelines no MongoDB.
      \end{itemize}
    \end{column}
  \end{columns}
\end{frame}

% ------------------------------------------------------------------------------
\begin{frame}{Arquitetura Integrada do Simulado: Cache-Aside}
  \vspace{-0.1cm}
  \centering
  \begin{tikzpicture}[
    box/.style={rectangle, draw=ScoutNavy, rounded corners=3pt, text centered, minimum height=1.0cm, line width=0.8pt, font=\scriptsize\bfseries},
    arrow/.style={thick, ->, >=stealth, color=ScoutGreen!85!black},
    lbl/.style={font=\tiny\bfseries, color=ScoutNavy!90, fill=white, inner sep=1.5pt}
  ]
    \node [box, fill=ScoutNavy!15, text width=2.8cm] (app) at (0, 0) {Aplicação Python\\(FastAPI / PyMongo)};
    \node [box, fill=ScoutGreen!20, text width=2.6cm] (redis) at (6.0, 0) {Redis 7.4\\(Cache em RAM / TTL)};
    \node [box, fill=ScoutGreen!35, text width=2.6cm] (mongo) at (6.0, -1.8) {MongoDB 8.0\\(Disco / BSON)};
    
    % Conexões App <-> Redis
    \draw [arrow] ([yshift=3pt]app.east) -- node[above, lbl] {1. Consulta Cache} ([yshift=3pt]redis.west);
    \draw [arrow, dashed] ([yshift=-3pt]redis.west) -- node[below, lbl] {Cache Hit ($<0.1\text{ ms}$)} ([yshift=-3pt]app.east);
    
    % Conexões App -> Mongo e Mongo -> Redis
    \draw [arrow] (app.south) |- node[near end, below, lbl] {2. Cache Miss: Lê Primário} (mongo.west);
    \draw [arrow] (mongo.north) -- node[right, lbl] {3. Salva no Cache com TTL} (redis.south);
  \end{tikzpicture}

  \vspace{0.25cm}
  \begin{columns}[t,onlytextwidth]
    \begin{column}{0.48\textwidth}
      \alert{Mecanismo de Consulta e Escrita}\par
      \footnotesize
      \begin{itemize}
        \item \textbf{Cache Hit:} Resposta instantânea em submilisegundo direto da RAM.
        \item \textbf{Cache Miss:} Busca indexada no MongoDB ($B\text{-Tree}$) e hidratação do Redis.
      \end{itemize}
    \end{column}
    \begin{column}{0.48\textwidth}
      \alert{Invalidação e Consistência}\par
      \footnotesize
      \begin{itemize}
        \item \textbf{Invalidação:} Deleção atômica da chave do Redis em alterações no MongoDB.
        \item \textbf{TTL Expirável:} Garante expiração automática contra dados obsoletos.
      \end{itemize}
    \end{column}
  \end{columns}
\end{frame}
